\section{Investigation Process}

Digital forensic investigations can be conceptually divided into a set of phases.\\ 
These phases do not represent a rigid or universal procedure, but rather a \textbf{conceptual framework} that helps investigators understand how to approach an investigation and what subtleties arise during each step.
\medskip

\noindent
Different standards propose slightly different models. Some frameworks describe four phases, others merge or rename certain steps.

\begin{quotation}
  \textbf{Author's npte:} \textit{ Porfessor Atzeny underlines that the goal is not to memorize a specific structure, but to understand the \textbf{underlying investigative logic}.}
\end{quotation}

\noindent
Well-known standards and guidelines include:

\begin{itemize}
\item \textbf{ACPO Guidelines} (United Kingdom)
\item \textbf{NIST Digital Forensics Framework}
\item \textbf{ISO standards} related to forensic processes
\end{itemize}

\noindent
Although these frameworks differ in terminology or number of steps, they describe broadly similar investigative processes.
\medskip

\noindent
The investigation process can be described using five phases:

\begin{enumerate}
\item Identification
\item Collection
\item Acquisition
\item Evaluation (or Examination)
\item Presentation
\end{enumerate}

\subsection{1. Identification}

The identification pphase consists of recognizing potential sources of digital evidence.\\
Typical sources include:

\begin{itemize}
\item storage devices (drives, removable media)
\item memory and running processes
\item mobile devices
\item network traffic
\item cloud storage and online services
\end{itemize}

\noindent
Investigatorsbegin forming preliminary hypotheses about relevance and also understanding the \textbf{volatility of information sources}. \\
Some evidence may disappear or change quickly, which means investigators must later prioritize acquisition accordingly.  

\medskip
\noindent
Whenever possible, investigators should gather contextual information \textbf{before interacting with the scene}, since direct interaction may alter digital evidence.
\\
Possible preliminary sources include:

\begin{itemize}
\item surveillance videos or photographs
\item social media content
\item organizational records or logs
\end{itemize}

\noindent
Another important factor is the evaluation of \textbf{technical knowledge of the suspect}.\\
If the suspect has limited technical expertise, the likelihood of sophisticated counter-forensic mechanisms may be lower.

\subsubsection{Scene Identification}

Once investigators reach the scene, the identification process continues with the recognition of actual devices and artifacts.\\
During this phase investigators may build a preliminary \textbf{graph of entities}, connecting devices, users, and network components that could be involved in the incident.

\subsubsection{Volatility and Prioritization}

Understanding volatility is essential because some evidence may disappear during normal system activity.
\\
Highly volatile sources should generally be acquired first.
\\
For instance:

\begin{itemize}
\item a powered-off device should normally remain powered off until acquisition
\item volatile memory may require specialized techniques to preserve its contents
\end{itemize}

\noindent
This prioritization begins during the identification phase and guides the subsequent collection process.

\subsubsection*{Example Scenario: Insider Data Exfiltration}

Consider an investigation inside a company where an employee is suspected of exfiltrating sensitive documents.\\
Potential evidence sources may include:

\begin{itemize}
\item employee workstations
\item USB storage devices
\item printers and shared devices
\item temporary files and browser caches
\item cloud storage services
\item internal file sharing systems
\end{itemize}

\noindent
Investigators may begin by collecting publicly available information about employees or suspects.\\  
Tools commonly used for this type of intelligence gathering include:

\begin{itemize}
\item \textbf{SpiderFoot}
\item \textbf{Maltego}
\item \textbf{Shodan}
\end{itemize}

\noindent
These tools help correlate information.
\medskip

\noindent
In managed environments investigators may also collect system information through administrative access.
\\
Typical diagnostic commands can reveal useful information about system activity, including:

\begin{itemize}
\item connected USB devices
\item mounted storage devices
\item open files and processes
\item active network connections
\end{itemize}

\noindent
For istance:

\begin{itemize}
\item \texttt{lsusb}, \texttt{mount}, \texttt{lsof}, \texttt{netstat} on Unix-like systems
\item administrative tools or PowerShell commands on Windows systems
\end{itemize}


\subsection{ 2. Collection Phase}

Once potential evidence sources have been identified, the next step is the \textbf{collection phase}.\\ 
The goal is to secure and isolate the identified evidence sources while minimizing any modification to their content.
\medskip
\noindent
A fundamental rule is to \textbf{avoid unnecessary interaction with the system}.

\subsubsection{Isolation of Evidence}

A central objective of the collection phase is the \textbf{isolation of evidence sources} to prevent further modification or destruction.\\
Isolation may occur at two levels:

\begin{itemize}
\item \textbf{Physical isolation}: preventing the suspect from accessing the device.
\item \textbf{Logical isolation}: preventing network communication with the device.
\end{itemize}

\noindent
Logical isolation may involve modifying network configurations such as firewall rules, network topology, or gateway settings to block communication with the device.

\subsubsection{Preservation of Evidence}

Investigators must also protect evidence from environmental or physical damage.\\
For examples:

\begin{itemize}
\item protecting magnetic storage devices from strong magnetic fields
\item handling fragile hardware carefully
\item preserving volatile memory when possible
\end{itemize}

In advanced scenarios, volatile memory may even be preserved using specialized techniques such as cryogenic preservation to delay data loss.

\subsubsection{Documentation and Evidence Tracking}

During the collection phase, all actions must be carefully documented as part of the chain of custody.\\

This documentation typically includes:

\begin{itemize}
\item who performed the action
\item when the action occurred
\item what operation was performed
\item which device was involved
\end{itemize}

\noindent
Devices should also be uniquely identified using information such as:

\begin{itemize}
\item model number
\item serial number
\item physical description
\end{itemize}

\subsubsection*{Forensic Equipment}

Forensic investigators typically use specialized equipment during the collection phase.\\
Examples include:

\begin{itemize}
\item hardware write blockers
\item forensic imaging devices
\item storage adapters for disk acquisition
\item signal jammers to prevent wireless communication
\end{itemize}


\subsection{ 3. Acquisition}

The \textbf{acquisition phase} consists of creating copies of the evidence collected during the previous phase.\\ 
These copies (artifacts) are used for analysis, allowing investigators to examine or manipulate data while preserving the original evidence.
\medskip

\noindent
The key requirement is that the replica must maintain \textbf{full integrity} with respect to the original data.

\subsubsection{Hashing and Integrity}

A central concept in this phase is \textbf{hashing}.  
Hash functions generate a digest that uniquely represents the content of a digital artifact.\\
During acquisition, hashing is typically performed:

\begin{itemize}
\item before the acquisition
\item after the acquisition
\item whenever the evidence is manipulated or transferred
\end{itemize}

\noindent
Even if it is possible to use different hash fuctions, the most commonly used in digital forensics are \textbf{MD5} and \textbf{SHA-family} (which is strictly racomanded).

\subsubsection{Static vs Live Acquisition}

A key decision during acquisition concerns whether the system should be analyzed in a \textbf{static} or \textbf{live} state.
\medskip

\noindent
\textbf{Static acquisition} involves powering off the device and copying the storage directly.\\  
This approach minimizes interaction with the system and is typically preferred when the relevant evidence resides on persistent storage.
\medskip

\noindent
However, static acquisition may not always be appropriate. If relevant data is located in volatile sources the \textbf{live acquisition} may be necessary.
\medskip

\noindent
In order to perform the acquisation phase, investigators use trust tools and.\\

\subsubsection{Documentation and Evidence Integrity}

All acquisition steps must be documented in detail, including:

\begin{itemize}
\item tools used
\item parameters and commands
\item timestamps
\item resulting hash values
\end{itemize}

\noindent
Throughout the acquisition phase, the \textbf{chain of custody} must be carefully maintained.\\
\underline{Every step in the handling of evidence must be traceable.}


\subsubsection{Example USB Acquisition}

Consider a simplified scenario where an employee is suspected of exfiltrating sensitive financial spreadsheets using a \textbf{USB flash drive}.\\
The investigation therefore focuses on acquiring the USB device in a forensically sound way.
\medskip

\noindent
\textbf{Preparation}\\
Before acquisition, the device must be clearly identified and documented:

\begin{itemize}
\item manufacturer, model, and serial number
\item photographs and physical description
\end{itemize}

\noindent
The investigator prepares:

\begin{itemize}
\item a trusted forensic workstation (Kali, CAINE)
\item a hardware write blocker
\item a dedicated forensic storage device
\end{itemize}

\noindent
\textbf{Acquisition Procedure}

\begin{enumerate}
\item Connect the USB drive through the \textbf{write blocker}
\item Create a \textbf{bit-by-bit forensic image} using a trusted tool (e.g., \texttt{dd}, \texttt{dc3dd}, FTK Imager)
\item Compute cryptographic hashes of the evidence and of the image
\item Verify that the hashes match
\item Document the entire procedure in the chain of custody
\end{enumerate}

\noindent
Example commands:

\begin{itemize}
\item Basic imaging with \texttt{dd}
\end{itemize}

\begin{verbatim}
dd if=/dev/sdb of=USB_image.dd
sha256sum USB_image.dd > USB_image.hash
\end{verbatim}

\begin{itemize}
\item Imaging with integrated hashing using \texttt{dc3dd}
\end{itemize}

\begin{verbatim}
dc3dd if=/dev/sdb of=USB_image.dd hash=sha256 log=USB_image.log
\end{verbatim}

\noindent
\textbf{Integrity Verification}\\
\noindent
After imaging, the hash of the image is compared with the original digest to confirm integrity.
\medskip

\noindent
\textbf{Documentation}\\
\noindent
The chain of custody must include:

\begin{itemize}
\item date, time, and location of acquisition
\item operator name and persons present
\item device identifiers (model, serial number)
\item tools, commands, and hash algorithms used
\end{itemize}

\subsection{4. Evaluation}

After acquisition, the investigation enters the \textbf{evaluation phase}, where the collected evidence is examined, correlated, and interpreted in order to reconstruct the events that occurred.

\subsubsection{Integrity Verification}

The first task is verifying whether the evidence is \textbf{authentic and admissible}.

\subsubsection{ Timeline Reconstruction}

The core objective of the evaluation phase is to reconstruct a coherent \textbf{timeline of events}.  
Investigators extract relevant information from multiple sources and combine them to identify:

\begin{itemize}
\item the sequence of actions,
\item the devices or actors involved,
\item the temporal relationships between events.
\end{itemize}

\noindent
Afeter these simple steps, investigators can begin to formulate hypotheses about what happened and why.

\subsubsection*{Anti-Forensic Challenges}

Modern systems introduce several obstacles to forensic analysis.  
Some mechanisms intentionally or unintentionally act as \textbf{anti-forensic factors}, including:

\begin{itemize}
\item disk encryption,
\item encrypted network communication (e.g.: HTTPS),
\item data obfuscation or deletion techniques.
\end{itemize}

\subsubsection*{Legal and Procedural Constraints}

The evaluation phase must comply also with \textbf{legal and procedural requirements}.
\\
Evidence may become inadmissible if:

\begin{itemize}
\item privacy regulations are violated,
\item investigative procedures exceed the legal mandate,
\item employees’ rights are breached during internal investigations.
\end{itemize}
\noindent
Therefore, forensic investigators must balance \textbf{technical analysis, evidentiary integrity, and legal compliance} throughout the evaluation process.


\subsection{5. Presentation}

The final phase of a forensic investigation is the \textbf{reporting and presentation of the results}.
\medskip

\noindent
The forensic analyst must adapt the report to the \textbf{target audience}. Different stakeholders may require different levels of detail:

\begin{itemize}
\item \textbf{technical experts}, who may require detailed procedures, tools, and commands used during the investigation;
\item \textbf{managers or executives}, who are primarily interested in the implications and consequences of the findings;
\item \textbf{judicial authorities}, who require a clear explanation of the evidence and its relevance within a legal framework.
\end{itemize}

\noindent
For this reason, multiple versions of the report may be necessary, each with a different level of abstraction and terminology.\\

\medskip
\noindent
Finally, the report itself becomes part of the \textbf{chain of custody}, so its integrity must be guaranteed.